% GENERATED FILE. Edit canonical lessons, then run npm run generate:books.
% canonical-source-hash: fnv1a64:5acd6854e09112d2
% canonical-lessons: PT-C355-vaga, PT-C355-curriculo, PT-C355-candidatura, PT-C355-candidato, PT-C355-contratar

\chapter{Procuro emprego}
\label{ch:pt-procuro-emprego}
\hlchaptermodality{355}

\begin{chapteropening}
\textbf{By the end of this chapter:} \emph{I can look for a job: find a vacancy, send a CV and an application, and talk about being hired.}

Complete the last lesson of chapter 355: I can look for a job: find a vacancy, send a CV and an application, and talk about being hired.
\end{chapteropening}

\section[a vaga]{a vaga — the vacancy, the opening; the wave (of heat)}

\label{lesson:PT-C355-vaga}



\begin{quote}
Before the new one: say the Portuguese for to ease, then the Portuguese for to make up for.
\end{quote}

\subsection*{You'll want to know: a vaga}
\textbf{a vaga} — \textquotedblleft{}the vacancy, the opening; the wave (of heat)\textquotedblright{}.

\begin{quote}
Há uma vaga para cozinheiro no restaurante da esquina.
\end{quote}

\begin{quote}
\emph{There is a vacancy for a cook at the restaurant on the corner.}
\end{quote}

Also \emph{uma vaga de calor} — a heatwave.

\begin{cousinweb}
\emph{vaga} \textquotedblleft{}vacancy\textquotedblright{} is built on the verb \emph{vagar} \textquotedblleft{}to fall vacant\textquotedblright{}, which comes from Latin \emph{vacāre} \textquotedblleft{}to be empty, to be free\textquotedblright{} (Lewis \& Short).
\end{cousinweb}

\begin{grammarlens}[title={What you've built}]
One more word for telling, explaining and describing at B1.
\end{grammarlens}

\subsection*{Your turn}
\begin{itemize}
\raggedright
  \item \emph{Say it:} \emph{a vaga}
  \item \emph{Say it:} \emph{a vaga}, once more
  \item \emph{Say it:} a sentence of your own with \emph{a vaga}
  \item \emph{Recall:} say the Portuguese for to ease, then the Portuguese for to make up for, then say \emph{a vaga} again
\end{itemize}

\begin{tcolorbox}[breakable,colback=teal!4,colframe=teal!35!black,title={Before you move on}]
What does a vaga mean? (\textquotedblleft{}The vacancy, the opening; the wave (of heat)\textquotedblright{}.) Say it once more.
\end{tcolorbox}
\section[o currículo]{o currículo — the CV, the résumé}

\label{lesson:PT-C355-curriculo}



\begin{quote}
Before the new one: say the Portuguese for to make up for, then the Portuguese for the vacancy.
\end{quote}

\subsection*{You'll want to know: o currículo}
\textbf{o currículo} — \textquotedblleft{}the CV, the résumé\textquotedblright{}.

\begin{quote}
Enviei o meu currículo a cinco empresas.
\end{quote}

\begin{quote}
\emph{I sent my CV to five companies.}
\end{quote}

Also \emph{o curriculum vitae} or \emph{o CV}. In schools, \emph{o currículo} is also the syllabus.

\begin{cousinweb}
\emph{currículo} is a learned borrowing from Latin \emph{curriculum} \textquotedblleft{}a running, a course, a race\textquotedblright{} (EtymDB, Lewis \& Short), built on \emph{currere} \textquotedblleft{}to run\textquotedblright{}: a life's course.
\end{cousinweb}

\begin{grammarlens}[title={What you've built}]
One more word for telling, explaining and describing at B1.
\end{grammarlens}

\subsection*{Your turn}
\begin{itemize}
\raggedright
  \item \emph{Say it:} \emph{o currículo}
  \item \emph{Say it:} \emph{o currículo}, once more
  \item \emph{Say it:} a sentence of your own with \emph{o currículo}
  \item \emph{Recall:} say the Portuguese for to make up for, then the Portuguese for the vacancy, then say \emph{o currículo} again
\end{itemize}

\begin{tcolorbox}[breakable,colback=teal!4,colframe=teal!35!black,title={Before you move on}]
What does o currículo mean? (\textquotedblleft{}The CV, the résumé\textquotedblright{}.) Say it once more.
\end{tcolorbox}
\section[a candidatura]{a candidatura — the application (for a job); the candidacy}

\label{lesson:PT-C355-candidatura}



\begin{quote}
Before the new one: say the Portuguese for the vacancy, then the Portuguese for the CV.
\end{quote}

\subsection*{You'll want to know: a candidatura}
\textbf{a candidatura} — \textquotedblleft{}the application (for a job); the candidacy\textquotedblright{}.

\begin{quote}
Fiz a candidatura ontem à noite.
\end{quote}

\begin{quote}
\emph{I put in the application last night.}
\end{quote}

\emph{Apresentar uma candidatura} — to submit an application.

\begin{cousinweb}
\emph{candidatura} is built on \emph{candidato} \textquotedblleft{}candidate\textquotedblright{}, from Latin \emph{candidātus} \textquotedblleft{}clothed in white\textquotedblright{} (Lewis \& Short), on \emph{candidus} \textquotedblleft{}shining white\textquotedblright{}.
\end{cousinweb}

\begin{grammarlens}[title={What you've built}]
One more word for telling, explaining and describing at B1.
\end{grammarlens}

\subsection*{Your turn}
\begin{itemize}
\raggedright
  \item \emph{Say it:} \emph{a candidatura}
  \item \emph{Say it:} \emph{a candidatura}, once more
  \item \emph{Say it:} a sentence of your own with \emph{a candidatura}
  \item \emph{Recall:} say the Portuguese for the vacancy, then the Portuguese for the CV, then say \emph{a candidatura} again
\end{itemize}

\begin{tcolorbox}[breakable,colback=teal!4,colframe=teal!35!black,title={Before you move on}]
What does a candidatura mean? (\textquotedblleft{}The application (for a job); the candidacy\textquotedblright{}.) Say it once more.
\end{tcolorbox}
\section[o candidato]{o candidato — the candidate, the applicant}

\label{lesson:PT-C355-candidato}



\begin{quote}
Before the new one: say the Portuguese for the CV, then the Portuguese for the application.
\end{quote}

\subsection*{You'll want to know: o candidato}
\textbf{o candidato} — \textquotedblleft{}the candidate, the applicant\textquotedblright{}.

\begin{quote}
Havia trinta candidatos para um só lugar.
\end{quote}

\begin{quote}
\emph{There were thirty applicants for a single place.}
\end{quote}

Feminine \emph{a candidata}.

\begin{cousinweb}
\emph{candidato} is a learned borrowing from Latin \emph{candidātus} \textquotedblleft{}clothed in white\textquotedblright{} (Lewis \& Short), from \emph{candidus} \textquotedblleft{}shining white\textquotedblright{}: Roman candidates for office wore a whitened toga.
\end{cousinweb}

\begin{grammarlens}[title={What you've built}]
One more word for telling, explaining and describing at B1.
\end{grammarlens}

\subsection*{Your turn}
\begin{itemize}
\raggedright
  \item \emph{Say it:} \emph{o candidato}
  \item \emph{Say it:} \emph{o candidato}, once more
  \item \emph{Say it:} a sentence of your own with \emph{o candidato}
  \item \emph{Recall:} say the Portuguese for the CV, then the Portuguese for the application, then say \emph{o candidato} again
\end{itemize}

\begin{tcolorbox}[breakable,colback=teal!4,colframe=teal!35!black,title={Before you move on}]
What does o candidato mean? (\textquotedblleft{}The candidate, the applicant\textquotedblright{}.) Say it once more.
\end{tcolorbox}
\section[contratar]{contratar — to hire, to take on}

\label{lesson:PT-C355-contratar}



\begin{quote}
Before the new one: say the Portuguese for the application, then the Portuguese for the candidate.
\end{quote}

\subsection*{You'll want to know: contratar}
\textbf{contratar} — \textquotedblleft{}to hire, to take on\textquotedblright{}.

\begin{quote}
A empresa contratou dez pessoas este mês.
\end{quote}

\begin{quote}
\emph{The company hired ten people this month.}
\end{quote}

\emph{Contratar alguém} — to hire someone. \emph{Um contrato} — a contract.

\begin{cousinweb}
\emph{contratar} is a learned borrowing from Latin \emph{contractāre} (EtymDB), a form of \emph{contrahere} \textquotedblleft{}to draw together, to make a contract\textquotedblright{}, built on \emph{trahere} \textquotedblleft{}to draw\textquotedblright{} (Lewis \& Short).
\end{cousinweb}

\begin{grammarlens}[title={What you've built}]
One more word for telling, explaining and describing at B1.
\end{grammarlens}

\subsection*{Your turn}
\begin{itemize}
\raggedright
  \item \emph{Say it:} \emph{contratar}
  \item \emph{Say it:} \emph{contratar}, once more
  \item \emph{Say it:} a sentence of your own with \emph{contratar}
  \item \emph{Recall:} say the Portuguese for the application, then the Portuguese for the candidate, then say \emph{contratar} again
\end{itemize}

\begin{tcolorbox}[breakable,colback=teal!4,colframe=teal!35!black,title={Before you move on}]
What does contratar mean? (\textquotedblleft{}To hire, to take on\textquotedblright{}.) Say it once more.
\end{tcolorbox}
